\chapter{Limits and Capital Allocation}\label{ch:fm:limits-and-capital-allocation}

A firm holds \$213 million of risk capital against six strategies, the expected shortfall of its annual P\&L in its worst one year in a hundred.
Measured one by one, the strategies would need \$394 million, and three of them would earn less than the firm's 15\% cost of capital. Measured by
what each contributes to the firm's worst years, two still do; a third clears it almost three times over; and the trend-following book, which
alone would need \$69 million, lowers the firm's capital by \$21 million, because it makes money in the years the others lose. Who is charged
for the tail decides which strategies the firm keeps.

\section{What a limit framework is for}

A risk limit (One Quant Book 6, chapter 29) caps one measure of one exposure. A firm has thousands of them; what holds them together is a
framework.

\begin{definition}[Limit framework]\label{def:fm:limits-and-capital-allocation:framework}
A firm's \emph{limit framework}\index{limit framework} is the set of its risk limits and the rules that relate them: the tree of levels (firm,
business, desk, strategy, trader) at which limits are set, the measures limited at each level, how a level's limits constrain those below it, who
may set, change and temporarily raise each limit, and what happens when one is breached.
\end{definition}

A framework allocates the firm's \omterm{def:fm:the-head-of-desks-job:appetite}{risk appetite} (chapter 6) down the risk hierarchy; its limits are the \omterm{def:fm:the-head-of-desks-job:appetite}{risk appetite statement} in numbers. It has
three jobs: to keep the firm inside its appetite, to make every desk's risk visible against an agreed number, and to force a conversation before a
limit is crossed, not after. A framework that only records breaches does the second job; one whose limits are raised whenever they are reached
does none (chapter 12).

\begin{omfigure}
\begin{tikzpicture}[font=\footnotesize,>=stealth,box/.style={draw,rounded corners,minimum height=0.62cm,align=center,fill=omDef!10}]
\node[box,fill=omProp!12,minimum width=3.6cm] (f) at (0,3) {firm\\{\scriptsize VaR, stress, liquidity}};
\node[box,minimum width=3.0cm] (d1) at (-3.2,1.6) {desk\\{\scriptsize VaR, stress, loss}};
\node[box,minimum width=3.0cm] (d2) at (3.2,1.6) {desk\\{\scriptsize VaR, stress, loss}};
\node[box,fill=gray!10,minimum width=2.4cm] (s1) at (-4.6,0.2) {strategy\\{\scriptsize VaR, loss, notional}};
\node[box,fill=gray!10,minimum width=2.4cm] (s2) at (-1.8,0.2) {strategy\\{\scriptsize sensitivities}};
\node[box,fill=gray!10,minimum width=2.4cm] (s3) at (1.8,0.2) {strategy\\{\scriptsize VaR, loss}};
\node[box,fill=gray!10,minimum width=2.4cm] (s4) at (4.6,0.2) {strategy\\{\scriptsize concentration}};
\draw[->] (f) -- (d1); \draw[->] (f) -- (d2); \draw[->] (d1) -- (s1); \draw[->] (d1) -- (s2); \draw[->] (d2) -- (s3); \draw[->] (d2) -- (s4);
\node[draw,rounded corners,fill=omThm!10,align=left,font=\scriptsize,anchor=west] at (-6.3,-1.3)
  {each limit: soft (warn) and hard (act) thresholds; temporary increases with an expiry;\\a child's exposure also counts against its parent's limit};
\end{tikzpicture}

\omcaption{A \omterm{def:fm:limits-and-capital-allocation:framework}{limit framework} as a tree: the firm's appetite is cascaded to desks and strategies, each level limiting the measures that fit its
risks. Model: \texttt{firm.limitalloc.Node}, \texttt{check}.}\label{fig:fm:limits-and-capital-allocation:tree}
\end{omfigure}

\section{Measures and levels: notional, sensitivities, VaR, stress, loss}

No single measure is enough, because each misses what another sees. Notional and position limits are simple and cannot be gamed by a model, but say
nothing about risk; sensitivity limits (delta, vega, DV01, credit spread sensitivity) are precise for small moves and blind to large ones; value at
risk and expected shortfall (One Quant Book 6, chapter 21) aggregate across risks but rest on a model and a history; stress losses (chapter 22 of
the same book) see scenarios the history has not; loss limits (One Quant Book 11, chapter 27) act only after the fact. A desk carries a few of
each, chosen for its risks.

\begin{definition}[Concentration limit]\label{def:fm:limits-and-capital-allocation:conc}
A \emph{concentration limit}\index{concentration limit} caps an exposure to one name, issuer, sector, counterparty or instrument, or a position
as a share of a market's volume or open interest, whatever the diversification elsewhere: it limits what could not be sold or hedged quickly.
\end{definition}

\begin{method}[Choosing a desk's limits]\label{met:fm:limits-and-capital-allocation:choose}
\begin{enumerate}
\item List the desk's risks and the moves that would hurt it most; choose one measure that sees each.
\item Set a VaR or ES limit for the aggregate, stress limits for the scenarios the firm's appetite names, sensitivity limits for the large
risks, a loss limit, and \omterm{def:fm:limits-and-capital-allocation:conc}{concentration limits} for what cannot be sold in days (chapter 28's cases show why).
\item Size each limit from the desk's share of the firm's appetite (chapter 6), and check that the desk can run its expected business at 60--80\%
of each.
\item Record who may change each limit and how temporary increases expire.
\end{enumerate}
\end{method}

\section{Utilisation, breaches and temporary increases}

\begin{definition}[Hard limit, soft limit, limit utilisation, limit breach]\label{def:fm:limits-and-capital-allocation:breach}
A \emph{hard limit}\index{hard limit} is a limit that must not be exceeded: an excess must be cured at once, by reducing the exposure, or
escalated for an approved increase. A \emph{soft limit}\index{soft limit} is a lower threshold whose crossing triggers a warning and a review, not
a mandatory action. \emph{Limit utilisation}\index{limit utilisation} is the exposure as a share of the limit. A \emph{limit breach}\index{limit
breach} is an exposure above a hard limit, recorded with its size, duration and resolution.
\end{definition}

The chapter's vol-selling strategy has a daily VaR limit of 8 with a soft threshold at 6. After a volatility spike its VaR reaches 7.6: a soft
breach and 95\% utilisation of the \omterm{def:fm:limits-and-capital-allocation:breach}{hard limit}. The desk asks for a temporary increase to 12 for sixteen days; utilisation falls to 63\% and the
position stays (\cref{fig:fm:limits-and-capital-allocation:var}). A temporary increase is a decision to take more risk for a stated time; one that
is renewed at each expiry is a permanent increase that nobody approved as one.

\begin{omfigure}
\begin{tikzpicture}
\begin{axis}[width=10.5cm,height=5.2cm,xlabel={\small day},ylabel={\small daily VaR (\$ million)},xmin=1,xmax=60,ymin=0,ymax=13,
  tick label style={font=\footnotesize},grid=major,grid style={gray!20},table/col sep=comma,
  legend cell align=left,legend style={font=\scriptsize,at={(0.5,-0.3)},anchor=north,/tikz/every even column/.append style={column sep=8pt}},legend columns=3]
\addplot[omDef,thick] table[x=day,y=var]{figdata/desk/07-limits-and-capital-allocation/varpath.csv};
\addplot[omProp,thick,dashed] table[x=day,y=soft]{figdata/desk/07-limits-and-capital-allocation/varpath.csv};
\addplot[omThm,thick,const plot] table[x=day,y=hard]{figdata/desk/07-limits-and-capital-allocation/varpath.csv};
\legend{daily VaR,soft limit,{hard limit (temporary increase days 25--40)}}
\end{axis}
\end{tikzpicture}

\omcaption{The vol-selling strategy's daily VaR through a volatility spike, against its \omterm{def:fm:limits-and-capital-allocation:breach}{soft limit} and its \omterm{def:fm:limits-and-capital-allocation:breach}{hard limit}, raised temporarily from 8
to 12 for days 25 to 40. Illustrative path; model \texttt{fm\_limits.var\_path}.}\label{fig:fm:limits-and-capital-allocation:var}
\end{omfigure}

The same check runs at every level. At the chapter's firm on the spike's first day, when vol selling's VaR has doubled to 6.4,
the six strategies' daily VaRs add up to \$21.5 million against
a firm limit of \$30 million (72\% utilisation), while vol selling alone is in soft breach: the firm is well inside its appetite and one desk is
not. A framework that checked only the firm would see nothing; one that checked only desks would never see the day when six desks at 70\%
add up to a firm breach.

\section{Allocating risk capital}

\begin{definition}[Economic capital]\label{def:fm:limits-and-capital-allocation:ec}
A firm's \emph{economic capital}\index{economic capital} is the capital it holds, by its own measure, against unexpected losses: typically the
expected shortfall or value at risk of its P\&L over a year at a high confidence level, computed on its own model of its risks.
\end{definition}

The chapter's firm (illustrative, \$ millions a year) runs six strategies: stat arb, trend following, credit carry, vol selling, macro and market
making, with expected P\&L of 11, 12, 8, 12, 6 and 10 and volatilities of 25, 30, 20, 22, 20 and 8, correlated at 0.2 except trend. In one year in
twenty a crash adds losses of 40, 60 and 110 to stat arb, credit carry and vol selling and a gain of 30 to trend. Its \omterm{def:fm:limits-and-capital-allocation:ec}{economic capital}, the 99\%
expected shortfall of annual P\&L over 200\,000 simulated years, is \$212.8 million for an expected P\&L of \$59.1 million.

Three allocations divide it (\cref{fig:fm:limits-and-capital-allocation:capital}). \emph{Stand-alone}: each strategy's own expected shortfall;
they add up to \$393.9 million, far more than the firm's, because the firm is diversified. \emph{Euler} (One Quant Book 6, chapter 20): each
strategy's expected P\&L in the firm's tail years, with the sign flipped; these add up exactly to the firm's. \emph{Incremental}: the firm's
expected shortfall minus what it would be without the strategy.

\begin{proposition}[Euler contributions of expected shortfall]\label{prop:fm:limits-and-capital-allocation:euler}
Let the firm's P\&L be $X=\sum_iX_i$ and $\mathrm{ES}_\alpha(X)=-\E[X\mid X\le q_{1-\alpha}(X)]$. For a continuous distribution, the Euler
contribution of strategy $i$, $w_i\,\partial\mathrm{ES}_\alpha(\sum_jw_jX_j)/\partial w_i$ at $w=\mathbf 1$, equals
$-\E[X_i\mid X\le q_{1-\alpha}(X)]$, and the contributions add up to $\mathrm{ES}_\alpha(X)$.
\end{proposition}

\begin{proof}[Partial proof]
Expected shortfall is positively homogeneous of degree one in the scales $w$, so Euler's theorem gives $\sum_iw_i\,\partial_i\mathrm{ES}=\mathrm{ES}$.
The derivative of the tail-conditional mean with respect to $w_i$ is the tail-conditional mean of $X_i$, because the quantile's own derivative
contributes nothing at the boundary (the density of $X$ at the quantile times a zero-length change of the tail set); the full argument, with the
regularity conditions, is Tasche's (One Quant Book 6, chapter 20). On scenarios, the contribution is the mean of $-X_i$ over the firm's worst
$(1-\alpha)n$ scenarios, and the sum is exactly the firm's scenario ES.
\end{proof}

\begin{omfigure}
\begin{tikzpicture}
\begin{axis}[width=11.5cm,height=5.4cm,ybar,bar width=7pt,ymin=-40,ymax=130,ylabel={\small capital (\$ million)},
  xtick={0,1,2,3,4,5},xticklabels={{stat arb},trend,{credit carry},{vol selling},macro,{market making}},
  xticklabel style={font=\scriptsize,text width=1.5cm,align=center},tick label style={font=\footnotesize},table/col sep=comma,
  enlarge x limits=0.1,grid=major,grid style={gray!20},legend cell align=left,
  legend style={font=\scriptsize,at={(0.5,-0.3)},anchor=north,/tikz/every even column/.append style={column sep=8pt}},legend columns=3]
\addplot[fill=gray!40,draw=gray] table[x=k,y=standalone]{figdata/desk/07-limits-and-capital-allocation/capital.csv};
\addplot[fill=omDef!55,draw=omDef] table[x=k,y=euler]{figdata/desk/07-limits-and-capital-allocation/capital.csv};
\addplot[fill=omThm!50,draw=omThm] table[x=k,y=incremental]{figdata/desk/07-limits-and-capital-allocation/capital.csv};
\draw[black] (axis cs:-0.6,0) -- (axis cs:5.6,0);
\legend{stand-alone ES,Euler contribution,incremental}
\end{axis}
\end{tikzpicture}

\omcaption{\omterm{def:fm:limits-and-capital-allocation:ec}{Economic capital} (99\% expected shortfall of annual P\&L) allocated to six strategies three ways. Stand-alone allocations add up to
\$394 million, the Euler contributions to the firm's \$213 million; trend following's contribution is negative because it gains in the firm's
worst years. Illustrative firm, 200\,000 simulated years. Data: \texttt{fm\_limits.allocation}.}\label{fig:fm:limits-and-capital-allocation:capital}
\end{omfigure}

\omcode{code/firm/limitalloc/firm_limitalloc.py}{83}{102}{Stand-alone, Euler and incremental allocations of expected shortfall on scenarios,
and the return on and charge for capital.}\label{lst:fm:limits-and-capital-allocation:alloc}

\section{Charging for capital}

\begin{definition}[Risk-adjusted return on capital, capital charge]\label{def:fm:limits-and-capital-allocation:raroc}
A strategy's \emph{risk-adjusted return on capital}\index{risk-adjusted return on capital} (RAROC) is its expected P\&L divided by the \omterm{def:fm:limits-and-capital-allocation:ec}{economic
capital} allocated to it. Its \emph{capital charge}\index{capital charge} is the hurdle rate (One Quant Book 6, chapter 19) times that capital,
deducted from its P\&L when its performance is measured and its pay pool sized.
\end{definition}

A strategy adds value when its RAROC exceeds the hurdle, that is when its expected P\&L exceeds its \omterm{def:fm:limits-and-capital-allocation:raroc}{capital charge}
(\cref{fig:fm:limits-and-capital-allocation:va}). With a 15\% hurdle:

\begin{center}
\small
\begin{tabular}{@{}lrrrrr@{}}
\toprule
& expected & \multicolumn{2}{c}{RAROC} & \multicolumn{2}{c}{value added}\\
strategy & P\&L & stand-alone & Euler & stand-alone & Euler\\
\midrule
stat arb & 10.9 & 16.7\% & 22.8\% & 1.12 & 3.74\\
trend & 12.0 & 17.4\% & (negative capital) & 1.66 & 15.24\\
credit carry & 8.0 & 10.5\% & 12.6\% & $-3.50$ & $-1.55$\\
vol selling & 12.1 & 9.8\% & 11.1\% & $-6.34$ & $-4.20$\\
macro & 6.0 & 12.6\% & 43.3\% & $-1.17$ & 3.93\\
market making & 10.0 & 84.9\% & (near zero capital) & 8.24 & 10.02\\
\bottomrule
\end{tabular}
\end{center}

\begin{omfigure}
\begin{tikzpicture}
\begin{axis}[width=11.5cm,height=5.2cm,ybar,bar width=9pt,ymin=-8,ymax=17,ylabel={\small value added (\$ million a year)},
  xtick={0,1,2,3,4,5},xticklabels={{stat arb},trend,{credit carry},{vol selling},macro,{market making}},
  xticklabel style={font=\scriptsize,text width=1.5cm,align=center},tick label style={font=\footnotesize},table/col sep=comma,
  enlarge x limits=0.1,grid=major,grid style={gray!20},legend cell align=left,
  legend style={font=\scriptsize,at={(0.5,-0.3)},anchor=north,/tikz/every even column/.append style={column sep=8pt}},legend columns=2]
\addplot[fill=gray!40,draw=gray] table[x=k,y=va_standalone]{figdata/desk/07-limits-and-capital-allocation/capital.csv};
\addplot[fill=omDef!55,draw=omDef] table[x=k,y=va_euler]{figdata/desk/07-limits-and-capital-allocation/capital.csv};
\draw[black] (axis cs:-0.6,0) -- (axis cs:5.6,0);
\legend{charged on stand-alone capital,charged on Euler capital}
\end{axis}
\end{tikzpicture}

\omcaption{Expected P\&L less a 15\% charge on allocated capital. On stand-alone capital three strategies destroy value; on Euler capital two do,
credit carry and vol selling, the two that lose most in the firm's worst years; macro and trend change sign or grow. Data:
\texttt{fm\_limits.value\_added}.}\label{fig:fm:limits-and-capital-allocation:va}
\end{omfigure}

Stand-alone charging penalises diversifiers and flatters strategies whose losses come together; Euler charging is consistent with the firm's actual
capital and rewards what reduces the firm's tail. It has its own trap: a strategy whose contribution is near zero or negative has an undefined or
negative RAROC and a \omterm{def:fm:limits-and-capital-allocation:raroc}{capital charge} that is a credit, and a desk paid on it will want more of it than the firm's tail can use (the contribution
changes as the book changes). Firms therefore cap credits, charge a floor, or use the incremental allocation for decisions about entering or
leaving a business and Euler for the running charge.

\begin{method}[Setting a capital charge]\label{met:fm:limits-and-capital-allocation:charge}
\begin{enumerate}
\item Compute the firm's \omterm{def:fm:limits-and-capital-allocation:ec}{economic capital} on scenarios that include its crash years, not only its recent history.
\item Allocate it by Euler contributions for the monthly charge; recompute the contributions when positions change materially.
\item Use incremental capital for decisions to add or close a strategy.
\item Charge capital at the hurdle rate in each desk's P\&L before its pay pool is sized (chapter 10), with a floor for strategies whose
contribution is near zero.
\end{enumerate}
\end{method}

A firm that re-sizes its strategies to maximise expected P\&L less a 15\% charge on its whole expected shortfall, each scaled between zero and
twice today's size (\texttt{firm.limitalloc.best\_scales}), grows four of them to the cap, cuts credit carry to 1.57 times and vol selling to
0.57 times, and raises the firm's value added from \$27.2 million to \$59.0 million a year. The cap does most of that work; the useful output is
the direction, the two tail-heavy books shrinking while the rest grow.

\section{Tutorial: six strategies and one tail}\label{tut:fm:limits-and-capital-allocation:tut}

\begin{tutorial}
\textbf{Goal.} Check a \omterm{def:fm:limits-and-capital-allocation:framework}{limit framework} on one day, allocate a firm's \omterm{def:fm:limits-and-capital-allocation:ec}{economic capital} three ways, and charge for it. \textbf{End state:} the table
of RAROCs and \cref{fig:fm:limits-and-capital-allocation:va}.
\begin{tutsteps}
\item \textbf{Limits.} \texttt{fm\_limits.framework()} builds the tree; \texttt{firm.limitalloc.check(framework(), daily\_var\_exposures(), 1)}
reports the soft breach and the firm's 72\% utilisation; \texttt{framework(temp=\dots)} adds a temporary increase of 4 on vol selling's
VaR limit until day 10.
\item \textbf{Scenarios.} \texttt{fm\_limits.scenarios()} simulates 200\,000 years of the six strategies with the crash.
\item \textbf{Allocation.} \texttt{allocation()} computes the three allocations and the RAROCs (\cref{lst:fm:limits-and-capital-allocation:alloc});
check that the Euler contributions add up to the firm's ES (\cref{prop:fm:limits-and-capital-allocation:euler}).
\item \textbf{Charge and scale.} \texttt{value\_added()} and \texttt{scaled()}.
\end{tutsteps}
\end{tutorial}

\textbf{What to change next.} Remove the crash and see the Euler allocation of trend turn positive; halve vol selling and recompute every
strategy's contribution.

\section{Build: limits and capital}\label{bld:fm:limits-and-capital-allocation:limitalloc}

\begin{build}
\bfield{Purpose} The firm's \omterm{def:fm:limits-and-capital-allocation:framework}{limit framework} as data and its risk capital as a price: which desk is inside its limits today, and what each strategy
costs in capital.
\bfield{Interface} \texttt{firm.limitalloc}: \texttt{Limit(measure, soft, hard)}, \texttt{Node(name, limits, children, temp)};
\texttt{check(node, exposures, day)}; \texttt{es(pnl, alpha)}, \texttt{allocate\_capital(scen, alpha)}; \texttt{raroc}, \texttt{capital\_charge};
\texttt{best\_scales(scen, alpha, h\_k, lo, hi)}; \texttt{to\_riskctl(node)}, the pre-trade part in the shape of One Quant Book 11's
\texttt{firm.riskctl} snapshot.
\bfield{Rules} A temporary increase applies only until its last day; a parent's exposure is the sum of its children's unless measured directly;
Euler contributions add up to the firm's ES exactly on the same scenarios.
\bfield{Acceptance tests} \texttt{code/firm/limitalloc/tests/}: soft and hard statuses, a parent breach, an expiring increase; Euler contributions
add up and stand-alone ones exceed the firm's ES; incremental allocations do not exceed stand-alone ones; the optimiser prefers the better
strategy; the export's shape.
\bfield{Stretch} Non-additive parent measures (a parent's VaR computed from the joint P\&L); an optimiser with the firm's limits as constraints;
a charge floor.
\end{build}

\begin{omsources}
\item E.~Zaik, J.~Walter, G.~Kelling and C.~James, ``RAROC at Bank of America: from theory to practice'', \emph{Journal of Applied Corporate
Finance} 9(2), 1996.
\item Financial Stability Board, \emph{Principles for an Effective Risk Appetite Framework}, 2013.
\item One Quant Book 6, chapters 19--21 and 29 (hurdle rate, Euler allocation, VaR and ES, the risk engine's limits).
\end{omsources}

\section{Exercises}

\begin{exercise}[$\star$]\label{exo:fm:limits-and-capital-allocation:1}
A strategy's daily VaR is 7.6 against a \omterm{def:fm:limits-and-capital-allocation:breach}{soft limit} of 6 and a \omterm{def:fm:limits-and-capital-allocation:breach}{hard limit} of 8. What is its utilisation, and what is its status? What if the \omterm{def:fm:limits-and-capital-allocation:breach}{hard
limit} is raised temporarily to 12?
\end{exercise}

\begin{exercise}[$\star$]\label{exo:fm:limits-and-capital-allocation:2}
Compute stat arb's RAROC on stand-alone and on Euler capital, and its value added at a 15\% hurdle on each.
\end{exercise}

\begin{exercise}[$\star$]\label{exo:fm:limits-and-capital-allocation:3}
The six strategies' stand-alone ES add up to \$393.9 million and the firm's is \$212.8 million. What is the diversification benefit, as a share of
the stand-alone sum?
\end{exercise}

\begin{exercise}[$\star\star$]\label{exo:fm:limits-and-capital-allocation:4}
At what hurdle rate would vol selling break even on its Euler capital? And credit carry?
\end{exercise}

\begin{exercise}[$\star\star$]\label{exo:fm:limits-and-capital-allocation:5}
Explain why trend following's Euler contribution is negative, and why a desk charged on it might take more risk than the firm wants.
\end{exercise}

\begin{exercise}[$\star\star$]\label{exo:fm:limits-and-capital-allocation:6}
Six desks each use 70\% of a VaR limit of 8 and the firm's limit is 30. Is the firm in breach if the desks' VaRs add? If they are correlated at
0.2 and each VaR is a multiple of its volatility?
\end{exercise}

\begin{exercise}[$\star\star\star$]\label{exo:fm:limits-and-capital-allocation:7}
\emph{Coding.} Rerun \texttt{fm\_limits.allocation} with no crash (\texttt{P\_CRASH = 0}). What are trend following's and vol selling's Euler
contributions now?
\end{exercise}

\begin{exercise}[$\star\star\star$]\label{exo:fm:limits-and-capital-allocation:8}
\emph{Find the flaw.} ``Every desk is inside its limits, so the firm is inside its \omterm{def:fm:the-head-of-desks-job:appetite}{risk appetite}.''
\end{exercise}

\section{Problem: Who Pays for the Tail}

\begin{problem}[{Weekend problem --- who pays for the tail}]\label{pb:fm:limits-and-capital-allocation:1}
The firm's head of risk must decide how to charge six strategies for \$213 million of \omterm{def:fm:limits-and-capital-allocation:ec}{economic capital}, and the vol-selling desk has just asked for
a temporary limit increase.

\medskip\noindent\textbf{Part I --- The framework.}
\begin{enumerate}
\item Define a \omterm{def:fm:limits-and-capital-allocation:framework}{limit framework}, and hard and \omterm{def:fm:limits-and-capital-allocation:breach}{soft limits}.
\item Give the vol-selling strategy's utilisation and status at the peak, before and after a temporary increase to 12.
\item Give the firm's VaR utilisation on the day of the spike.
\item Why must a framework check both the firm and each desk?
\item What turns a temporary increase into an unapproved permanent one?
\end{enumerate}

\medskip\noindent\textbf{Part II --- Capital.}
\begin{enumerate}[resume]
\item Define \omterm{def:fm:limits-and-capital-allocation:ec}{economic capital} and give the firm's, with its expected P\&L.
\item Give each strategy's stand-alone, Euler and incremental capital.
\item State and prove \cref{prop:fm:limits-and-capital-allocation:euler}; check that the contributions add up.
\item Why does the stand-alone sum exceed the firm's capital?
\item Which strategy's capital changes sign between the stand-alone and Euler allocations, and why?
\end{enumerate}

\medskip\noindent\textbf{Part III --- The charge.}
\begin{enumerate}[resume]
\item Define RAROC and a \omterm{def:fm:limits-and-capital-allocation:raroc}{capital charge}.
\item Which strategies destroy value at a 15\% hurdle under each allocation?
\item Give each strategy's value added under both.
\item What goes wrong with a near-zero or negative Euler contribution, and what do firms do about it?
\item Which allocation should decide whether to close a strategy?
\end{enumerate}

\medskip\noindent\textbf{Part IV --- The decision.}
\begin{enumerate}[resume]
\item Give the best scales in $[0,2]$ and the value added before and after.
\item Which part of that result would you trust, and which not?
\item Should the vol-selling desk get its temporary increase?
\item State the \emph{named result}: the hurdle rates at which credit carry and vol selling break even on stand-alone and on Euler capital, and
trend's capital under both.
\item In two sentences, write the head of risk's recommendation.
\end{enumerate}
\end{problem}

\section{Interview questions}

\begin{interviewq}[$\star$ \iqroles{risk}]\label{iq:fm:limits-and-capital-allocation:1}
What is the difference between a soft and a \omterm{def:fm:limits-and-capital-allocation:breach}{hard limit}? Give an example of each at a trading desk.
\end{interviewq}

\begin{interviewq}[$\star$ \iqroles{trader, risk}]\label{iq:fm:limits-and-capital-allocation:2}
Why does a desk carry several kinds of limit instead of one VaR limit?
\end{interviewq}

\begin{interviewq}[$\star\star$ \iqroles{risk, researcher}]\label{iq:fm:limits-and-capital-allocation:3}
Three strategies each have a stand-alone ES of 50 and the firm's ES is 90. How would you allocate the 90, and why not 50 each?
\end{interviewq}

\begin{interviewq}[$\star\star$ \iqroles{bank, risk}]\label{iq:fm:limits-and-capital-allocation:4}
A strategy earns 12 a year on Euler capital of 108. At a 15\% hurdle, does it add value? What would change your answer?
\end{interviewq}

\begin{interviewq}[$\star\star$ \iqroles{researcher}]\label{iq:fm:limits-and-capital-allocation:5}
Why can a strategy's Euler contribution to expected shortfall be negative? Is that a reason to grow it without limit?
\end{interviewq}

\begin{interviewq}[$\star\star\star$ \iqroles{risk}]\label{iq:fm:limits-and-capital-allocation:6}
How would you estimate each strategy's contribution to a firm's tail when the tail is driven by rare crashes that the last five years do not
contain?
\end{interviewq}
